% DRAFT — Chapter 2 body expanded per THESIS_EXPANSION_PLAN.md (17 Aug 2026).
% Replaces lines 1–517 of chapters/02_paper1.tex on approval; the appendix is untouched.

\chapter{Measuring School Culture: Warmth and Strictness in English Secondary Schools}
\label{ch:paper1}
\resetcounterformats

\begin{center}
  \textit{Damian Phelan%
    \thanks{Department of Economics, University College London.
      Primary supervisor: Michela Tincani. Secondary supervisor: Lars Nesheim.
      [\textit{INSERT: acknowledgements, funding statement.}]}}
\end{center}

\begin{chapterabstract}
School culture is widely held to matter for what pupils learn, but it is
rarely measured. This paper constructs measures of two dimensions of the
culture of English secondary schools, warmth and strictness, and asks how far
each can be recovered from data that exist for every school. For \NTierOne{}
schools a research visit and a headteacher interview yield a gold standard in
two parts, an enacted score from observation and an espoused score from the
head's own account; inspection reports, behaviour policies and websites are
scored for the national population by language-model instruments, and
prediction models fitted where both exist carry the gold standard to the
rest. Three findings follow. Strictness is legible in public text and warmth
is not: inspection-report strictness is correlated with observed strictness at
$r = \CorrOfstedStrictness$ and a text-only model predicts it out of sample at
$\RidgeBStrictness$, while no text source is correlated with observed warmth.
Espoused and enacted culture diverge, and more for warmth: headteachers rate
their schools \GapWarmth{} points warmer on a ten-point scale than observers
record and \GapStrictness{} points stricter, and every public source tracks the
espoused warmth measure rather than the enacted one. And across every version
of every language-model instrument tested, agreement with reference labels was
unrelated to correlation with the observed measure.
\end{chapterabstract}

% ---------------------------------------------------------------
\section{Introduction}
\label{sec:p1_intro}

School culture is widely held to matter for what pupils learn, but it is
rarely measured. This paper constructs measures of two dimensions of the
culture of English secondary schools, warmth and strictness, for the full
population of around 3,330 state secondary schools, and asks how far each
dimension can be recovered from data that exist for every school.

Warmth is the relational quality of a school: whether pupils are known,
cared for and supported by the adults around them. In a lesson it shows in
whether the teacher uses pupils' names, how praise is given and how the
teacher and the class speak to one another; between lessons it shows in how
staff and pupils deal with each other in corridors and at break. Strictness is
the school's behavioural structure: whether expectations of conduct are clear,
enforced and met. In a lesson it shows in how much misbehaviour is tolerated
before a sanction, how much disruption occurs and how pupils respond when they
are corrected; outside lessons it shows in the order of arrival, movement and
the dining hall. Strictness in this sense is a matter of clarity and
consistency rather than severity: a school that runs centralised detentions,
a phone ban and silent corridors, applied by every member of staff, is strict
whether or not it is also warm.

The two dimensions are the axes of the parenting typology of
\citet{baumrind1971} and \citet{maccoby1983}, which classifies parenting by
demandingness and responsiveness and finds that the combination of both, the
authoritative style, produces the best child outcomes across cognitive and
non-cognitive domains; \citet{majumder2015} and \citet{zhang2020} find the
same pattern in educational attainment, with each dimension carrying an
independent effect. The transfer to schools is direct: demandingness is
strictness and responsiveness is warmth. \citet{pellerin2005} and
\citet{gregory2010,gregory2011} apply the typology to schools and find that
pupils in schools combining support with structure report less victimisation
and more engagement, and \citet{gottfredson2005} reviews the related evidence
on school discipline. What makes the two dimensions measurable in a school,
as opposed to a family, is that both are enacted in public. A researcher in a
lesson can see whether names are used and how disruption is met, and a
researcher in a corridor can see whether pupils move calmly and whether staff
speak to them; neither requires access to anything private. Warmth and
strictness are treated here as distinct dimensions rather than the ends of one
axis, because a school can be warm without being orderly and orderly without
being warm, and whether the two co-vary is a question the data answer rather
than a premise. A third dimension, the quality of teaching practice, is
observed alongside them; it is not part of culture as defined here but is the
obvious rival explanation for any association between culture and results,
and it is carried for that reason.

The difficulty is that none of this is recorded in any administrative
dataset. What is recorded, for every school, is text: an inspection report
written by an external observer, a behaviour policy written by the school, and
a website written by the school for parents. What is recorded for a smaller
set of schools, because it was collected for this study, is direct observation
and a structured interview with the headteacher. The design follows from that
asymmetry. For \NTierOne{} schools there is a research visit, comprising
structured lesson and out-of-lesson observations, and a headteacher interview;
from these a gold-standard measure of each dimension is constructed. Because
the interview and the visit disagree substantially about the same school, the
gold standard is two measures, not one: an \emph{enacted} score, from what an
observer recorded, and an \emph{espoused} score, from what the headteacher
reported. For a further \NTierTwo{} schools there is the interview alone, and
for the remaining national population only the public text. For those schools
the gold standard is estimated from what is observable, using prediction
models fitted on the schools where both the predictors and the gold standard
exist, and the paper reports how well the predictions reproduce the gold
standard out of sample.

Three findings follow. First, the two dimensions are recoverable to very
different degrees. Strictness scored from inspection-report text is
correlated with observed strictness at $r = \CorrOfstedStrictness$, and a
text-only prediction model retains a leave-one-out correlation of
$\RidgeBStrictness$ with the gold standard. Warmth is not recoverable from
any public source: no text-derived warmth score is correlated with observed
warmth, and the text-only model has no out-of-sample skill (leave-one-out
$R^2$ of $-0.02$). Second, the enacted and espoused measures diverge in a
systematic and dimension-specific way. Headteachers rate their schools nearly
two points warmer on a ten-point scale than observers record and about one
point stricter; the interview and the visit are correlated at
$\GoldWarmthSplit$ for warmth and $\GoldStrictnessSplit$ for strictness.
Every public source, including a survey of parents, tracks the espoused warmth
measure and not the enacted one. Third, the language-model instruments used to
score the text carry a measurement result of their own: an instrument's
agreement with its reference labels and its correlation with the observed
measure are unrelated, and the best-agreeing version of an instrument was
usually the worse measure.

The paper sits between five literatures and adds something to each. The
school-climate literature measures culture through pupil, teacher and parent
surveys \citep{hoy2001,cohen2009,thapa2013}. Those surveys have produced
consistent associations between climate and wellbeing, attendance and
attainment, but they capture perceptions rather than practice, are open to
social desirability, and are costly at national scale; \citet{vanhoutte2005}
argues that climate, the aggregate of members' perceptions, should be kept
apart from culture, the shared norms that govern behaviour, and it is culture
in that sense that the measures here aim at. Administrative proxies such as
exclusion rates \citep{machin2020} conflate culture with intake, and the
inspection grade folds many dimensions into one verdict. This paper's
contribution to that literature is a measure of enacted culture built from
observation rather than perception, with the perception data, from
headteachers and from parents, used as a comparison against it rather than as
a component of it.

The school-effectiveness tradition that begins with \citet{coleman1966} and
\citet{rutter1979} located the difference schools make not in resources but
in what Rutter called ethos, the expectations and routine practice a school
sustains, and later syntheses converged on a list that includes an orderly
environment and a supportive climate \citep{edmonds1979,sammons1995,teddlie2000}.
That strand has stayed largely qualitative because ethos has been measured
by case study; the contribution here is to make two of its correlates
measurable across a national population. Structured classroom observation,
in particular the CLASS framework of \citet{hamre2007}, measures emotional
support and classroom organisation directly and predicts pupil gains at
secondary level \citep{allen2013} and in English primary schools
\citep{sammons2008}; the observation schedules used here draw on those two
domains, but aggregate across lessons and out-of-lesson settings to
characterise a school rather than a teacher, and to the author's knowledge no
study has combined such observation with a leadership interview and public
text to separate what a school does from what it says. The literature on
high-expectation, high-structure charter and academy models
\citep{abdulkadiroglu2011,angrist2010,dobbie2013,fryer2014} finds large
effects of a bundle of practices in which strict behaviour norms travel with
warm, individual relationships; it is the origin of the strictness dimension
and the reason to expect it to matter, but it has not measured the two
dimensions apart. And the text-as-data literature in economics
\citep{gentzkow2019} supplies the tools for reading the public record at
scale, with the standing warning that general-purpose methods misread
domain-specific text \citep{loughran2011,muddiman2019}; this paper applies
large language models rather than dictionary or topic methods, tests every
instrument against direct observation, and reports what that choice costs and
buys in measurement terms. Relative to all five, the contribution is a measure
of enacted culture with a national text-based extension whose validity is
tested against observation rather than assumed, together with a set of school
level scores that are the explanatory variables of the following chapter.

The paper proceeds as follows. Section~\ref{sec:p1_data} describes the data.
Section~\ref{sec:p1_measurement} constructs the gold standard, the text-based
instruments and the prediction models, and reports how well the predictions
reproduce the gold standard. Section~\ref{sec:p1_results} reports what the
measures show. Section~\ref{sec:p1_conclusion} concludes.

% ---------------------------------------------------------------
\section{Data}
\label{sec:p1_data}

The paper uses three kinds of data on English state secondary schools: a
research visit and a headteacher interview conducted for this study, three
corpora of public text, and the Department for Education's school-level
administrative panel. \Cref{tab:tier_summary} sets out coverage.

\textbf{Visits and interviews.} Between November 2023 and April 2025 I
interviewed the headteacher, or where unavailable the senior leader
responsible for behaviour and culture, at 304 schools, and visited
\NTierOne{} of them. Schools were recruited by direct approach to
headteachers by email and telephone and by referral from participating
schools; no quota was applied by region, school type or inspection grade.
Three-quarters of interviews and half of visits took place in the 2024/25
academic year. Every headteacher interviewed was asked at the end of the
interview whether the school would host a visit, so the visited schools are a
self-selected subset of the interviewed ones.

A visit lasted a school day and was made by a team of two, occasionally
three, researchers. The team observed lessons across the timetable, typically
four to six in a day and five at the median, spread across Years 7 to 11 and
across subjects, each lesson observed for its full length by every researcher
present rating independently on a fixed schedule; the team also observed
arrival, movement between lessons, break and lunch, and completed a second
schedule once for the day. The lesson schedule has nine items bearing on
warmth and strictness (for example the frequency with which the teacher uses
pupils' names; the level of misbehaviour tolerated before a first sanction)
and nine on teaching practice (questioning, feedback, explanation, structure
and the like); the out-of-lesson schedule has eight items covering
transitions, breaks, corridors, arrival and the dining hall. Every item is
rated 1 to 5 against written anchors at 1, 3 and 5, and item ratings are
averaged across researchers and across lessons to give one value per school.
The interview lasted about an hour and followed a fixed schedule of eleven
open questions on the school's approach to behaviour, relationships, wellbeing
and teaching, ten statements on which the headteacher rated their own school
on a five-point agreement scale (for example ``Sanction systems are applied
consistently across the school''; ``Pupils in this school feel that their
teachers care about them''), and a checklist of ten named behaviour systems
(centralised detentions, an internal exclusion room, a mobile phone ban and
so on) recorded as present or absent. Where a second researcher attended, both
scored the open questions independently. The full schedules are reproduced in
\cref{app:2A:interview_guide,app:2A:visit_protocol}, and inter-rater agreement
on the observation and interview items is reported with the measures in
Section~\ref{sec:p1_measurement}.

\textbf{Public text.} Three corpora cover the national population. Ofsted
inspection reports were obtained for 3,288 of the 3,312 open schools
(99.3 per cent); the remainder opened after 2023 and had not been inspected.
Behaviour policy documents were obtained for 3,363 schools, one document per
school, located by crawling from the school's website and selected by a rule
that prefers the current whole-school behaviour policy over annexes,
predecessor versions and sixth-form-only documents (\cref{app:2A:crawl});
about one document in ten is a trust-wide policy adopted verbatim by the
school. School websites were crawled for 3,370 schools, taking the home,
values and welcome pages, which are the pages on which a school describes
itself to prospective parents; 3,284 crawls contain scorable prose, and where a
crawl returned no prose the school is flagged rather than dropped, so that
missingness is a variable and not a silent exclusion.

\textbf{Administrative data.} All schools are matched on their unique
reference number to a school-year panel built from the Department for
Education's published tables for 2015/16 to 2024/25, carrying pupil
composition, prior attainment, Progress~8 and its components, workforce and
finance. This paper uses the panel only to describe the samples; the
following chapter uses it for outcomes and controls.

\textbf{Representativeness.} \Cref{tab:2A:representativeness} compares the
visited and interviewed schools with the national population. Both samples
resemble the population on school type, phase, admissions, size, deprivation
and urban location, and both over-represent London (45 per cent of visited
schools against 15 per cent nationally), which shows in a higher share of
pupils with English as an additional language. Both also lean toward stronger
schools: 25 per cent of visited schools held an Outstanding inspection grade
as at August 2024 against 15 per cent nationally, and their mean Progress~8
in 2018/19 was $+0.17$ against $+0.04$. That lean is the expected result of
recruiting through headteachers, since confident schools say yes, and it has
two consequences. Estimates from the visited sample apply most directly to a
population somewhat stronger than the national average; and because the
sample under-represents the weakest schools, the range of both culture and
outcomes is restricted, which works against finding associations rather than
for them.

\Cref{tab:score_dist} reports descriptive statistics for the measures
constructed in Section~\ref{sec:p1_measurement}, for the visited schools.

\begin{table}[htbp]
\centering
\small
\begin{tabular}{lcccccc}
\toprule
Group & $N$ & FSM\% & EAL\% & Median size & Outstanding\% & Good\% \\
\midrule
Visited and interviewed & 103 & 29.2 & 27.8 & 1{,}118 & 25 & 65 \\
Interviewed only & 200 & 27.2 & 17.9 & 1{,}052 & 15 & 75 \\
Rest of the national population & 3{,}028 & 30.1 & 18.2 & 1{,}050 & 14 & 70 \\
\midrule
All schools & 3{,}332 & 29.9 & 18.4 & 1{,}051 & 15 & 70 \\
\bottomrule
\end{tabular}
\caption{Distribution of schools across the three groups and selected observable characteristics (2023--24 school year). The groups are mutually exclusive: the first received both a headteacher interview and a researcher visit, the second an interview only, and the third is the rest of the national population. FSM\%: mean proportion eligible for free school meals.
EAL\%: mean proportion with English as an additional language. Size: median
number of pupils on roll. Outstanding\%/Good\%: percentage rated Outstanding or
Good at most recent inspection, among schools with a grade.}
\label{tab:tier_summary}
\smallskip
\begin{minipage}{0.92\linewidth}
\footnotesize\textit{Notes:} The groups sum to one fewer than the total: one interviewed school whose head left the statement battery blank belongs to the total but to no group.
\end{minipage}
\end{table}


\begin{table}[htbp]
\centering
\caption{Distribution of sub-scores and of the enacted and espoused scores for
Tier~1 full-data schools ($n = 103$).
Sub-scores are on a $[0,5]$ scale; the enacted and espoused scores on a $[0,10]$ scale.}
\label{tab:score_dist}
\footnotesize\setlength{\tabcolsep}{4pt}
\begin{tabular}{llcccccc}
\toprule
Score & Description & $N$ & Mean & SD & Min & Median & Max \\
\midrule
$W1$ & In-lesson staff--pupil warmth & 103 & 3.24 & 0.45 & 1.92 & 3.22 & 4.47 \\
$W2$ & Warmth in transitions and breaks & 103 & 3.32 & 0.68 & 1.83 & 3.33 & 4.67 \\
$W3$ & HT espoused warmth (statements) & 103 & 4.20 & 0.47 & 2.33 & 4.33 & 5.00 \\
$S1$ & In-lesson behaviour management & 103 & 3.69 & 0.49 & 2.12 & 3.70 & 4.75 \\
$S2$ & Out-of-lesson behaviour management & 103 & 3.33 & 0.50 & 2.10 & 3.30 & 4.60 \\
$S3$ & Behaviour management systems (count) & 103 & 3.11 & 0.73 & 0.62 & 3.12 & 4.38 \\
$S4$ & HT espoused strictness (statements) & 103 & 3.95 & 0.58 & 1.75 & 4.00 & 5.00 \\
$T1$ & In-lesson teaching quality & 103 & 3.48 & 0.42 & 2.49 & 3.51 & 4.51 \\
$SC$ & HT espoused staff climate (statements) & 103 & 3.91 & 0.78 & 1.50 & 4.00 & 5.00 \\
\midrule
Warmth & \emph{enacted} & 103 & 6.56 & 0.99 & 4.33 & 6.60 & 8.81 \\
Strictness & \emph{enacted} & 103 & 7.02 & 0.87 & 4.85 & 7.04 & 8.96 \\
Teaching & \emph{enacted} & 103 & 6.97 & 0.83 & 4.98 & 7.02 & 9.02 \\
Warmth & \emph{espoused} & 103 & 8.39 & 0.95 & 4.67 & 8.67 & 10.00 \\
Strictness & \emph{espoused} & 103 & 7.90 & 1.16 & 3.50 & 8.00 & 10.00 \\
Staff climate & \emph{espoused} & 103 & 7.81 & 1.55 & 3.00 & 8.00 & 10.00 \\
\bottomrule
\end{tabular}
\begin{minipage}{\textwidth}\vspace{0.5em}\footnotesize
\textit{Notes.} The enacted scores are built from the visit sub-scores
($W1$, $W2$, $S1$, $S2$, $T1$) and the espoused scores from the interview
statement sub-scores ($W3$, $S4$; $SC$ is the staff-climate score), as set
out in \cref{eq:enacted,eq:espoused}. The systems count $S3$ enters no
score. No school contributes a sub-score to both sets.
\end{minipage}
\end{table}


% ---------------------------------------------------------------
\section{Measurement}
\label{sec:p1_measurement}

This section builds the measures in the order the design requires. It first
constructs the gold-standard scores from the visits and interviews and reports
their reliability. It then sets out how the public text is scored by
language-model instruments and how those instruments were tested. It ends
with the prediction models that carry the gold standard to schools that were
not visited, and the evidence on how well they do so.

\subsection{The gold standard}

Every observation item and interview statement is rated 1 to 5. Sub-scores
are item means; the dimension scores are sub-score means, rescaled to 0--10.
Two sets are constructed and kept apart. The \emph{enacted} scores use the
observations alone: enacted warmth averages in-lesson warmth (names, praise,
quality of interaction) and out-of-lesson warmth (interactions during
transitions and breaks); enacted strictness averages in-lesson strictness
(misbehaviour tolerated, disruption, response to sanction, respect) and
out-of-lesson strictness (sanction use, corridors, arrival, dining hall,
recreational space); enacted teaching averages the nine lesson items on
practice. The \emph{espoused} scores use the headteacher's statement ratings
alone: espoused warmth averages the three warmth statements (rewards applied
consistently; pupils feel safe; pupils feel their teachers care), espoused
strictness the four strictness statements (sanctions applied consistently;
behaviour in lessons; behaviour outside lessons; teachers feel respected). The
two remaining scored statements, on staff turnover and morale, form a
staff-climate score; the tenth statement, on parental engagement, enters no
score, and the interview yields no espoused measure of teaching. Item
definitions and the mapping to sub-scores are in \cref{app:2A:guide}.

To make the construction concrete, take the in-lesson strictness sub-score
of one school. In each of five lessons two researchers rated the level of
misbehaviour tolerated before a sanction, the frequency of disruption, the
pupils' response to correction and how respectfully they treated the teacher.
Misbehaviour and disruption are reverse-scored so that a higher value is
stricter, and where no sanction was needed in a lesson the response item is
filled from the disruption item, since a calm lesson in which no correction
was required is evidence about the school and not a missing value. The
twenty ratings per researcher are averaged to a lesson value, the lesson
values across lessons, and the researchers' values across researchers, giving
one in-lesson strictness figure on the 1--5 scale. That figure is averaged
with the out-of-lesson strictness figure and doubled to give the school's
enacted strictness on 0--10. Espoused strictness is simpler: the
headteacher's ratings on the four strictness statements are averaged and
doubled. Every school that received a visit has an enacted score; every
school that gave an interview has an espoused score; \NTierOne{} schools
have both.

Two components of the interview are recorded but excluded from the scores. The
question-level ratings of how fully the headteacher answered each open
question measure articulacy rather than culture: they load on one general
factor that tracks the interviewer's rating of the head's confidence
($r = 0.34$ to $0.47$) and correlate with the observed scores at zero. A count
of the ten named behaviour systems correlates with observed strictness at
$-0.01$, against $0.41$ for the statement battery, and is likewise excluded:
what systems a school names is not how strictly it runs. Both are retained in
the data as variables in their own right.

Inter-rater agreement is substantial. Two researchers rated every lesson and
every out-of-lesson period independently; pairwise Pearson correlations across
the observation items range from $0.47$ to $0.81$, and within-one-point
agreement is close to 90 per cent throughout
(\cref{tab:irr_classroom,tab:irr_outside}). Agreement is highest on the items
with the most concrete anchors, the frequency of disruption ($0.79$) and the
level of misbehaviour tolerated ($0.70$), and lowest on the teaching-practice
items that ask the observer for the most interpretation, such as the use of
resources and the choice of methods ($0.53$ and $0.55$). The interview
items were rated by the interviewer and an observer, with correlations of
$0.61$ to $0.74$ and within-one-point agreement above 90 per cent
(\cref{tab:irr_interview}). Enacted warmth and strictness are correlated at
$r = \CorrWarmthStrictnessEnacted$ across the \NTierOne{} visited schools,
and the espoused pair at $\CorrWarmthStrictnessEspoused$ across the same
schools; part of each correlation is common-method variance, since both
scores come from the same observers on the same day, so the separation of
the two dimensions is if anything understated.

The gold standard is two measures rather than one because the interview and
the visit disagree about the same school. Enacted and espoused warmth are
correlated at $r = \GoldWarmthSplit$ ($p = \GoldWarmthSplitP$) and enacted and
espoused strictness at $\GoldStrictnessSplit$ ($p = \GoldStrictnessSplitP$),
on the \NTierOne{} schools with both. Two measures sharing at most a sixth of
their variance are not fallible readings of one quantity. A weighted average
is the right operation for two noisy measures of a common target, because
their errors partly cancel; it is the wrong operation for two measures of
different targets, because it produces a third quantity that is neither,
whose composition depends on the weight chosen. The two are therefore carried
separately throughout, and their divergence is itself a result reported in
Section~\ref{sec:p1_results}. Two schools show what the divergence looks like
in the ratings. At the first, the headteacher rated the school 8.0 on warmth
and 8.0 on strictness; the observers, across their lessons and the day,
recorded 8.0 and 8.5, and the in-lesson and out-of-lesson sub-scores agree
with one another to within half a point. At the second, the headteacher rated
warmth 9.3 and strictness 7.8, near the top of the scale, and the observers
recorded 4.7 and 5.9: in-lesson warmth averaged 2.2 out of 5 and out-of-lesson
warmth 2.5, while the head's ratings on the three warmth statements averaged
4.7 out of 5. Nothing in the interview distinguishes the two heads. Across the
\NTierOne{} schools, four in five have a warmth gap larger than one point on
the 0--10 scale and two in five a gap larger than two. Both measures are also
correlated with an external administrative record where one exists: the
espoused staff-climate score is correlated with the school's recorded teacher
turnover at $r = -0.52$, and at $-0.48$ conditional on the full control set,
which is evidence that headteachers report accurately on what they know.

\subsection{Scoring the public text}

For schools without a visit the gold standard must be predicted from what is
observable, and the richest observable is text. Each of the three corpora is
scored by a language-model instrument that reads one document about one
school and returns a 1--5 score for warmth and for strictness (and, for
inspection reports, teaching). All scoring uses gpt-4o-mini at zero
sampling temperature; every score records the prompt version, a hash of the
rubric text and the model version, and identical inputs are served from
cache, so that any score in the dataset can be traced to the exact instrument
that produced it.

The pipeline is the same for every corpus. A document is retrieved and
reduced to plain text; for inspection reports every stated grade and every
field that would let the model infer the grade is stripped from the body,
so that the score is read from the prose and not from the verdict. The text
is presented to the model with the instrument's rubric, and the model's answer
is converted to a band. Two architectures are used at that last step, and the
choice between them is empirical. \emph{Prose} instruments give the model a
written marking scheme and take its band; the inspection-report and
behaviour-policy instruments are of this kind, and each scores its constructs
in a single call on the full document. \emph{Decomposed} instruments never ask
for a band: the model answers a fixed set of factual questions, each answer
supported by a verbatim quotation verified against the document, and the band
is assigned by a fixed rule in code. The website strictness instrument, for
example, asks whether the page names a specific consequence for a specific
behaviour, whether it describes a rule at the level of a particular (a phone,
a corridor, a uniform item), whether it uses the vocabulary of expectation,
and so on; each answer must be backed by a quotation that appears in the
page, unbacked answers are discarded, and the band is assigned by a fixed
ladder over the answers that survive. The website and interview-transcript instruments are of this kind.
The operative rubric and question set of every instrument is reproduced in
\cref{app:2A:prompts}.

Instruments were tuned against reference labels and then tested three ways.
The reference labels were produced by three independent instances of a
larger language model (Claude, Opus class), each given only the written rule
and the documents, blind to one another and to every model score, with the
majority label as the target; the written rules were drafted and revised by
the author, who ruled on the boundary cases the raters could not and
labelled the development packs personally. Agreement figures below are
therefore model--model agreement, and the final arbiter of every instrument
is not agreement but the gold standard. The author's agreement with his own
labels of the same documents two weeks apart is $\kappa_w = 0.77$, which
bounds what any rater can achieve against them at about $0.88$. Tuning and
testing were kept apart by design: each instrument was developed on a
labelled pack of documents, every version of it was archived, and every
archived version was then scored once on a fresh pack that no version had
been tuned on.

The first test is agreement out of sample: weighted kappa against the
reference labels on documents no version of the instrument was tuned on.
Every in-sample figure had read high, and the out-of-sample figures were
lower for every instrument; \cref{tab:instruments_adopted} reports them for
the instruments retained.\footnote{About 85 archived versions across eleven
instrument families were rescored in this way.} The second is validity:
correlation with the enacted gold standard on the visited schools, and, for
one construct, with an external register. The website instrument for the
prominence of religious faith recovers the official register of religious
character with an area under the ROC curve of 0.960, which establishes that
the method reads what is in the text; where a construct fails against its
criterion, the failure is a property of the text. The third is a check that no
field available to the model carries the inspection verdict: report bodies
are stripped of every stated grade before scoring, and how far the verdict
can still be read from the prose that remains is reported below.

\begin{table}[htbp]
\centering
\small
\caption{The scoring instruments adopted for the national tier. Agreement is
weighted kappa against the reference labels, measured out of sample on
documents no version was tuned on. Validity is each instrument's correlation
with the enacted gold standard on the visited schools, or, where an official
register exists, agreement with it.}
\label{tab:instruments_adopted}
\begin{tabular}{llclc}
\toprule
\textit{Source} & \textit{Construct} & \textit{Architecture} &
\textit{Validity} & $\kappa_w$ \\
\midrule
Inspection report & Strictness & Prose & visits $r = +0.40$ & 0.64 \\
Inspection report & Warmth & Prose & visits $r = -0.08$ & 0.45 \\
Inspection report & Teaching & Prose & visits $r = +0.06$ & 0.75 \\
\addlinespace[3pt]
Behaviour policy & Strictness & Prose & visits $r = -0.08$ & 0.76 \\
Behaviour policy & Warmth & Prose & visits $r = +0.01$ & 0.76 \\
\addlinespace[3pt]
Website & Strictness & Decomposed & visits $r = +0.09$ & 0.66 \\
Website & Warmth & Decomposed & visits $r = -0.04$ & 0.69 \\
Website & Religious character & Prose classifier & register, 97\% agree & --- \\
\addlinespace[3pt]
Interview transcript & Strictness & Decomposed & --- & 0.63 \\
Interview transcript & Warmth & Decomposed & --- & 0.56 \\
\bottomrule
\end{tabular}
\begin{minipage}{0.92\linewidth}
\vspace{4pt}\scriptsize
\textit{Notes}: ``Prose'': the model reads a written marking scheme and
returns a band in one call. ``Decomposed'': the model answers factual
questions with verified quotations and the band is assigned by a rule in
code. Kappas are model--model agreement against the majority label of three
blind model raters. The inspection-report warmth score also tracks the
inspection grade ($r = -0.40$) from grade-stripped text, and is used as a
description of the report rather than a predictor. The website warmth
instrument does not separate its two lowest bands reliably. The interview
transcript instruments are methodological comparisons only; the espoused
scores of record come from the statement battery, not the transcripts.
\end{minipage}
\end{table}


The testing produced four results that governed which instruments were kept,
summarised in \cref{tab:instrument_findings}, and each is a finding about
language-model measurement rather than about schools.

The first is that agreement with labels and correlation with the gold
standard are unrelated across versions, and the best-agreeing version was
usually the worse measure. The inspection-report strictness instrument that
agreed best with its labels correlated with observed strictness at $0.07$;
the prose version, which agreed less well, correlated at $0.43$. The reason
is that a rubric can be tuned toward what labellers reliably notice in a
document, and what labellers reliably notice is not always what a school
does; agreement measures the first, and only the visits measure the second.
Every instrument here was therefore selected on its correlation with the
observed measure, with agreement treated as a floor and never as the
objective.

The second is that instruments which ask the model for facts outperform those
that ask for judgements. Where a question can be answered by pointing at a
sentence, whether the page names a consequence, whether the report says
pupils arrive on time, the model's answers are stable, verifiable and
reproducible; where a question asks how strict or how warm the document
sounds, the answer carries the model's own generosity, and that generosity
moves between versions and between models. Counting verified facts in code
and assigning the band from the count removes the judgement from the model
and puts it in a rule that can be read.

The third is that scoring related constructs in separate calls can produce
a correlation between them that the gold standard does not show. Enacted
warmth and strictness are correlated at $\CorrWarmthStrictnessEnacted$; the
single-call prose instrument, which reads a document once and returns both
scores, gives $0.54$; the separate-call instruments return $0.08$ to $0.22$.
A model that reads for warmth in one call and strictness in another has no
sense of the school as a whole, and returns two readings of two passages
rather than one reading of one school; the single-call figure sits closer to
the observed one for that reason.

The fourth is that the model is part of the instrument. An unchanged rubric
moved between models fell from $\kappa_w = 0.63$ to $0.12$, from substantial
agreement to near none, with the text and the labels held fixed. Instruments
are therefore pinned to the model they were tested with, and a change of
model is treated as a new instrument to be tested again rather than as an
upgrade.

Two limitations of the retained instruments are stated once here and apply
wherever their scores are used. The inspection-report warmth score tracks the
inspection grade at $-0.40$ from text that never states it, where human
raters of the same text track it at $-0.25$; it is carried as a descriptive
measure and never as a clean predictor. The website warmth instrument leaves
its bottom two bands nearly indistinguishable at the model's limit.

\begin{table}[htbp]
\centering
\small
\caption{Four results from testing the language-model instruments, each with
the evidence that established it. All figures are out of sample.}
\label{tab:instrument_findings}
\begin{tabular}{p{4.2cm}p{9.2cm}}
\toprule
\textit{Result} & \textit{Evidence} \\
\midrule
Agreement with labels and validity are unrelated &
Inspection-report strictness: the version with the highest label agreement
correlates with observed strictness at $0.07$; the prose version at $0.43$.
Across eleven instrument families the best-agreeing version was usually the
worse measure. \\
\addlinespace[3pt]
Facts outperform judgements &
The one instrument family whose out-of-sample agreement improved across
versions did so when judgement questions were replaced by factual ones.
Judgement questions track the model's own generosity rather than the
document. \\
\addlinespace[3pt]
Separate calls can manufacture independence &
Enacted warmth and strictness correlate at $0.59$ in the visits. The
single-call prose instrument gives $0.54$; instruments scoring
each construct in a separate call return $0.08$ to $0.22$. \\
\addlinespace[3pt]
The model is part of the instrument &
An unchanged rubric moved from one scoring model to another fell from
$\kappa_w = 0.63$ to $0.12$. Every instrument is pinned to the model it was
tested with; one (faith prominence) is kept on an older model for that reason. \\
\bottomrule
\end{tabular}
\end{table}


\subsection{Predicting the gold standard}

For interviewed schools without a visit (Model~A) and for the national
population (Model~B), the enacted score of each dimension is predicted by
ridge regression fitted on the \RidgeTrainN{} visited schools with complete
predictors. Model~A's predictors are the espoused sub-scores, the systems
count and the two text scores for the dimension; Model~B's are the two text
scores alone. Ridge is used rather than least squares because the predictors
are few, correlated and measured with error, and a small penalty on the
coefficients trades a little bias for a useful reduction in variance; the
penalty is chosen by generalised cross-validation. Predictions are evaluated by
leave-one-out cross-validation on the same schools, which is the honest
metric here: each school's prediction comes from a model that never saw it,
so the reported correlation is what the model would achieve on a new school
(\cref{tab:interview_model_loo,tab:national_model_loo}). The fitted
coefficients are in \cref{app:2A:ridge_spec}.

Strictness is predictable and warmth is not. Model~A predicts enacted
strictness at a leave-one-out $r = \RidgeAStrictness$ and Model~B, from text
alone, at $\RidgeBStrictness$. For warmth Model~A reaches $\RidgeAWarmth$ and
Model~B has no out-of-sample skill: its leave-one-out $R^2$ is $-0.02$, no
better than assigning every school the sample mean, and its predictions
collapse toward a constant with a standard deviation under a tenth of the
gold standard's. Among the predictors, the espoused strictness score and the
inspection-report strictness score do roughly equal work ($0.24$ and $0.23$
per standard deviation), which is to say that a headteacher's account of
whether the school's systems are applied consistently and an inspector's
description of conduct carry about the same information about what an
observer will find; the systems count contributes nothing.

\begin{table}[htbp]
\centering
\caption{Leave-one-out cross-validated performance of the interview-tier Ridge
calibration model (Model~A) for warmth and strictness. Features: interview
sub-scores ($W3^*$, $S3$, $S4$, $\omega$) plus Ofsted and behaviour-policy LLM
holistic scores (6 features per dimension). Training and evaluation sample: $N = 100$. Of the 102 full-data schools, 2 schools lack a scored Ofsted report and are dropped.}
\label{tab:interview_model_loo}
\small
\begin{tabular}{lcccr}
\toprule
Dimension & LOO-CV $r$ & LOO-CV $R^2$ & RMSE & $\lambda$ \\
\midrule
Warmth     & 0.767 & 0.588 & 0.573 & 10 \\
Strictness & 0.701 & 0.489 & 0.492 & 1 \\
\bottomrule
\end{tabular}
\end{table}

\begin{table}[htbp]
\centering
\caption{Leave-one-out cross-validated performance of the national-tier Ridge
calibration model (Model~B) for warmth and strictness. Features: Ofsted LLM
score, behaviour-policy LLM score and website LLM score (3 features per
dimension). Target: the \emph{enacted} score from the school visit; no interview
input enters either side. Training and evaluation sample: $N = 101$. Of the 103 full-data schools, 2 schools lack a scored Ofsted report and are dropped.}
\label{tab:national_model_loo}
\small
\begin{tabular}{lcccr}
\toprule
Dimension & LOO-CV $r$ & LOO-CV $R^2$ & RMSE & $\lambda$ \\
\midrule
Warmth     & -0.720 & -0.023 & 0.937 & 1000 \\
Strictness & 0.319 & 0.099 & 0.803 & 10 \\
\bottomrule
\end{tabular}
\end{table}

Because Model~B's warmth predictions carry no signal and Model~A's carry
little, neither model's predictions are used as regressors in the following
chapter. That chapter uses the enacted scores directly for the visited
schools and the raw inspection-report strictness score for the national
extension. The prediction models are best read as the measurement result they
are: publicly available text can rank English secondary schools on strictness
and cannot rank them on warmth.

% ---------------------------------------------------------------
\section{Results}
\label{sec:p1_results}

This section reports what the measures show. Three results stand out: the
enacted and espoused measures diverge, and they diverge more for warmth than
for strictness; every public source of text tracks the espoused warmth measure
and none tracks the enacted one; and behaviour policy text is, if anything,
inversely related to observed practice.

\textbf{Distributions.} \Cref{tab:score_dist} reports the measures for the
visited schools and \cref{fig:score_distributions} the distributions of the
espoused and inspection-report scores at scale. Enacted warmth across
the \NTierOne{} visited schools has mean 6.6 and standard deviation 1.0 on the
0--10 scale; enacted strictness has mean 7.0 and standard deviation 0.9.
Espoused scores across the \NInterviewed{} interviewed schools sit higher
(means 8.5 and 8.0) with standard deviations of 0.9 and 1.0. The two enacted
dimensions are correlated at $r = \CorrWarmthStrictnessEnacted$: warm schools
tend also to be orderly ones, which is consistent with the authoritative-school
framework and inconsistent with a single warm-to-strict axis, and the four
quadrants of the joint distribution are all populated. Between-school
variation is genuinely narrow in the English system, and score compression
partly reflects that: about two-thirds of secondary schools are graded Good
and a further sixth Outstanding, and common inspection, curriculum and
training pressures compress the distribution of any culture measure. A
standard deviation of enacted warmth is therefore about one point on the
ten-point scale, which is the unit in which the gaps below should be read.

\begin{figure}[htbp]
  \centering
  \includegraphics[width=\textwidth]{fig_A_score_distributions.pdf}
  \caption{Distributions of the espoused warmth and strictness scores across
    the \NInterviewed{} interviewed schools (top panels, 0--10 scale) and of
    the inspection-report warmth and strictness scores across the national
    population (bottom panels, 1--5 scale). Dashed lines mark sample means.}
  \label{fig:score_distributions}
\end{figure}

\textbf{Enacted and espoused culture diverge.} On the \NTierOne{} schools
with both, headteachers rate their schools warmer than observers record by
$\GapWarmth$ points on the 0--10 scale and stricter by $\GapStrictness$ (paired
$p < 0.001$ for both). A gap of $\GapWarmth$ points is close to two standard
deviations of the enacted measure: the typical headteacher places their school
where an observer would place a school near the top of the visited range. The
two sources are correlated at $\GoldWarmthSplit$ for warmth and
$\GoldStrictnessSplit$ for strictness. The gap is therefore larger, and the
agreement lower, for warmth: what a headteacher says about relationships is a
poorer guide to what a visitor finds than what they say about behaviour.
Within the interview, the items on which heads report reliably are those about
the consistency of practice (sanctions applied consistently, $r = 0.39$ with
observed strictness; rewards applied consistently, $0.29$), not those about
how people feel (pupils feel safe, $0.16$; teachers care, $0.11$, neither
significant). Headteachers report accurately on whether their systems are
applied and inaccurately on relational climate, which is the pattern one
would expect if the first is something a head can check and the second is
something a head must infer.

\textbf{Public text tracks talk, not practice.} \Cref{tab:source_criterion}
sets each source's warmth and strictness score against the enacted and the
espoused criterion. The pattern is uniform. For strictness, the
inspection-report score is correlated with the observed measure at
$\CorrOfstedStrictness$ and with the espoused measure at
$\CorrOfstedStrictEspoused$; behaviour policies and websites carry no
strictness signal against either. For warmth, no source is correlated with
the observed measure, and the inspection report and the parent survey are
correlated with the espoused one ($0.29$ and $0.31$). The parent survey is
the fourth independent respondent group, after observers, headteachers and
inspectors, and it shows the same crossover: parents' reports of their
children being happy and safe track what heads say about warmth ($0.31$) more
than what observers record ($0.16$), and track observed strictness ($0.28$)
better than observed warmth. \Cref{fig:crossover_warmth} draws the warmth half
of the table. What is written and reported about a school's warmth, by
inspectors, by parents and by the school itself, tracks the school's account
of itself and not the classroom. One reading is that the school's account is
the common source: a head who describes the school as warm describes it that
way to inspectors and to parents, and both repeat it. Another is that
inspectors and parents see the same public face that the interview elicits,
and that the classroom and the corridor are where warmth is enacted or not.
The data do not separate the two, and both imply that a national measure of
enacted warmth cannot be read from any of these sources.

\begin{table}[htbp]
\centering
\small
\caption{Each text source's warmth and strictness score correlated with the
enacted (observed) and espoused (headteacher-reported) gold-standard measure.
Enacted correlations use the visited schools; espoused correlations use the
interviewed schools. The parent survey has no strictness item; its warmth
composite is also shown against enacted strictness.}
\label{tab:source_criterion}
\begin{tabular}{lcccc}
\toprule
 & \multicolumn{2}{c}{\textit{Warmth score vs}} & \multicolumn{2}{c}{\textit{Strictness score vs}} \\
\cmidrule(lr){2-3}\cmidrule(lr){4-5}
\textit{Source} & enacted & espoused & enacted & espoused \\
 & $(n \approx 101)$ & $(n \approx 299)$ & $(n \approx 101)$ & $(n \approx 299)$ \\
\midrule
Ofsted inspection report & $-0.07$ & $+0.29$ & $+0.40$ & $+0.38$ \\
Behaviour policy         & $-0.08$ & $+0.00$ & $+0.01$ & $+0.00$ \\
School website           & $-0.03$ & $-0.07$ & $+0.09$ & $+0.07$ \\
Parent survey            & $+0.16$ & $+0.31$ & \multicolumn{2}{c}{($+0.28$ vs enacted strictness)} \\
\bottomrule
\end{tabular}
\begin{minipage}{0.86\linewidth}
\vspace{4pt}\scriptsize
\textit{Notes}: Pearson correlations. Text scores are the instruments of
record in \cref{tab:instruments_adopted}. Enacted and espoused measures are
defined in Section~3. Correlations of $|r| \geq 0.20$ are significant at the
five per cent level at these sample sizes.
\end{minipage}
\end{table}


\textbf{Behaviour policies describe aspiration.} The behaviour-policy warmth
score is uncorrelated with observed warmth ($-0.08$; $-0.02$ on the revised
instrument used in the following chapter) and negatively correlated with the
parent measure ($-0.04$, $p = 0.03$); the strictness score is uncorrelated
with observed strictness ($0.01$ and $-0.08$ respectively). One school in ten
publishes a policy authored by its trust rather than by the school, and a
template shared across a trust's schools cannot vary with the culture of each. Written behaviour policies are governance documents that state
intent, and schools that state the warmest intent are not, on these measures,
the warmest schools. The same result appears in the following chapter as a
small negative association between behaviour-policy warmth and pupil
progress.

\textbf{What the text does and does not contain.} Two further results
qualify the inspection-report scores. The warmth score tracks the inspection
grade at $-0.40$ from report text stripped of every stated grade, where human
raters of the same text track it at $-0.25$: inspectors write about
relationships in a register that moves with their overall verdict, and the
model reads that register. The score is therefore a characterisation of the
report rather than a clean measure of the school. The strictness score also
tracks the verdict ($-0.49$) but, unlike warmth, it tracks observed practice
as well, and in the following chapter its association with pupil progress
survives conditioning on the grade. Finally, the one website construct with an
external criterion, the prominence of religious faith, is recovered from
website text with an area under the ROC curve of 0.960: the method reads what
is there, and warmth is not.

\begin{figure}[htbp]
  \centering
  \includegraphics[width=0.85\textwidth]{fig_crossover_warmth.pdf}
  \caption{Each source's warmth score correlated with the espoused criterion
  (filled) and the enacted criterion (open). Every source tracks what
  headteachers report better than what observers record.}
  \label{fig:crossover_warmth}
\end{figure}

% ---------------------------------------------------------------
\section{Conclusion}
\label{sec:p1_conclusion}

This paper constructs measures of warmth and strictness for English secondary
schools from direct observation, headteacher interview and public text, and
asks how far each dimension can be recovered from data that exist for every
school. Three findings summarise it.

Strictness is legible in public text and warmth is not. An inspection
report's strictness score is correlated with observed strictness at
$\CorrOfstedStrictness$, and a text-only model predicts it out of sample at
$\RidgeBStrictness$; no text source is correlated with observed warmth, and
the text-only warmth model has no out-of-sample skill. The asymmetry is a
fact about the sources, not the instruments: inspectors record conduct and
enforcement in detail and describe relationships in a uniform positive
register, schools' own documents describe intent, and websites are not written
about the thing observers see. For the school-climate literature the result
is a caution: the sources that are cheap at national scale carry one of the
two dimensions and not the other, and a survey of parents carries the same
one.

Espoused and enacted culture diverge, and more for warmth. Headteachers
over-state warmth relative to observation by nearly two points and agree with
observers only weakly on it, while their reports of behaviour systems being
applied consistently track observed strictness. Every public source, including
a survey of parents, tracks the espoused warmth measure and not the enacted
one. The gap is dimension-specific and appears in four independent kinds of
evidence, which suggests it is a feature of how schools present themselves
rather than an artefact of any one instrument. For the school-effectiveness
tradition, which has measured ethos largely through leaders' and teachers'
accounts, the result says which parts of those accounts can be trusted: the
parts about whether systems are applied, and not the parts about how people
feel.

The instruments carry a measurement result of their own. Across every version
of every language-model instrument tested, agreement with reference labels
was unrelated to correlation with the observed measure, and the best-agreeing
version was usually the worse measure; instruments that ask for facts
outperformed those that ask for judgements; scoring related constructs in
separate calls produced correlations between them that the observations do
not show; and an unchanged rubric moved between models lost most of its
agreement. These results apply to any study that scores text with a language
model against human labels, and they argue for testing every such instrument
against an external criterion rather than against agreement alone.

The measures have limits that the following chapter inherits. The visited
sample is \NTierOne{} schools, weighted toward London and toward stronger
schools; the enacted scores come from one visit each, and a school on a
different day might have been rated differently; the inspection-report scores
are read from prose that carries the inspection verdict; and the reference
labels for the text instruments were produced by a language model rather than
by human raters, with the observed measures as the final check. Within those
limits, the enacted scores are a reliable, two-dimensional measure of school
culture for the visited schools, and the inspection-report strictness score
extends one of the two dimensions to the national population. The following
chapter asks what the two dimensions, so measured, are worth in pupil
progress.
